\textbf{(CAP-COLUNI 2023 -- Q17)} No livro \textit{O Andar do Bêbado}, o autor Leonard Mlodinow relata uma dificuldade que os gregos tinham para escrever os números. Eles usavam uma espécie de código alfabético no qual as primeiras nove das 24 letras maiúsculas do alfabeto grego representavam os números que chamamos de 1, 2, ..., 9. As seguintes nove letras representavam os números que chamamos de 10, 20, ..., 90. E, as seis últimas letras, representavam as centenas 100, 200, ..., 600.

Abaixo, estão representadas, ordenadamente da esquerda para a direita, estas 24 letras:
\[
A\ B\ \Gamma\ \Delta\ E\ Z\ H\ \Theta\ I\ K\ \Lambda\ M\ N\ \Xi\ O\ \Pi\ P\ \Sigma\ T\ \Upsilon\ \Phi\ X\ \Psi\ \Omega.
\]

De acordo com o texto, é CORRETO afirmar que o resultado da subtração entre $\Omega\Xi H$ e $XM\Theta$ é:

\begin{enumerate}
	\item $\Upsilon K I$
	\item $\Phi\Lambda\Delta$
	\item $\Upsilon\Lambda I$
	\item $\Phi P E$
\end{enumerate}
